% !TEX program = pdflatex
% MLCAD Workshop Paper — IEEE Conference Format
\documentclass[conference]{IEEEtran}
\usepackage{cite}
\usepackage{amsmath,amssymb,amsfonts}
\usepackage{graphicx}
\usepackage{textcomp}
\usepackage{xcolor}
\usepackage{booktabs}
\usepackage{multirow}
\usepackage{hyperref}
\usepackage{balance}

\graphicspath{{./}}

\begin{document}

\title{Instrumented ML-Driven PPA Optimization:\\Pre-Registered Experiments, Negative Results,\\and Cross-Version EDA Generalization}

\author{
\IEEEauthorblockN{Oguzhan Tekin}
\IEEEauthorblockA{Independent Researcher\\
Toronto, ON, Canada\\
\textit{oguzhantekin@gmail.com}}
}

\maketitle

\begin{abstract}
We present an instrumented framework for ML-driven Power, Performance, and Area (PPA)
optimization in ASIC design that pre-registers experimental outcomes before execution,
separates surrogate model roles by architectural fit, and validates predictions through
cross-version EDA probes. Using OpenROAD Flow Scripts (ORFS) on ASAP7 with four RTL
designs, the full experimental arc—from surrogate training through Bayesian optimization
to EDA verification—cost \$8.35 in cloud compute.

Six findings emerge, each with pre-registered success criteria: (1)~GBDT$\rightarrow$RL
and GAT$\rightarrow$BO role separation yields an 85.1-point IQM improvement (from
$-$17.6 to +67.5) over the original monolithic architecture; (2)~graph-topology-only GAT achieves LODO-CV MAE of
0.153, matching criticality-augmented variants; (3)~criticality edge weights provide no
benefit beyond graph topology, with the mechanism being timing-repair-induced
node-count--clock correlation ($\rho{=}{-}1.0$); (4)~surrogate-only Bayesian optimization with
qLogNEHVI proposes 4 and 3 new Pareto candidates for AES and JPEG
respectively, passing the surrogate's dominance criterion; (5)~ORFS version changes (v3.0$\rightarrow$26Q1) produce erratic PPA
divergence including timing-sign flips, a methodological hazard for ML-for-EDA
reproducibility; (6)~all 7 BO-proposed candidates lie on the Pareto front under a
different ORFS version than training, with trade-off directions matching surrogate
predictions, and one candidate strictly dominating a control configuration.

Four pre-flight sanity checks caught four distinct bugs before they could propagate
incorrect results. We argue that the methodology of pre-registered, instrumented
experimentation is itself a contribution to ML-for-EDA practice.
\end{abstract}

\begin{IEEEkeywords}
PPA optimization, surrogate modeling, Bayesian optimization, graph neural networks,
OpenROAD, ASAP7, reproducibility, pre-registration
\end{IEEEkeywords}

%======================================================================
\section{Introduction}
%======================================================================

Machine learning for electronic design automation (ML-for-EDA) has produced numerous
results showing that surrogate models can predict PPA metrics and guide design-space
exploration~\cite{kahng2021,openroad2019,cheng2021}. However, the field faces
reproducibility challenges: published work rarely controls for EDA tool version
sensitivity, and negative results are infrequently reported.

We present a framework that addresses these gaps through three methodological
commitments:

\begin{enumerate}
\item \textbf{Pre-registration.} Each experimental phase has success criteria committed
to a plan document before execution. Outcomes are classified as pass/fail against
pre-specified thresholds.

\item \textbf{Negative-result reporting.} Findings that contradict initial hypotheses
(criticality redundancy, ORFS version sensitivity) are reported with the same rigor
as positive findings, including mechanistic explanations.

\item \textbf{Instrumented sanity checking.} Every phase begins with a pre-flight
diagnostic that must pass before scaling. Four of four sanity checks caught distinct
bugs that would have produced plausible-but-wrong results.
\end{enumerate}

The framework uses two surrogate architectures—Gradient-Boosted Decision Trees (GBDT)
and Graph Attention Networks (GAT)—assigned to different optimization roles based on
their architectural strengths: GBDT for fast, tabular-input reinforcement learning
reward evaluation, and GAT for uncertainty-aware Bayesian optimization over
graph-structured netlists.

The full experimental arc, from data generation through EDA verification, ran on AWS
spot and on-demand instances for a total compute cost of \$8.35.

%======================================================================
\section{Related Work}
%======================================================================

\textbf{Surrogate-based PPA prediction.}
GNNs operating on circuit netlists have shown promise for PPA
prediction~\cite{ustun2020,huang2021,wu2022,chhabria2022}. OpenABC-D~\cite{chowdhury2021}
provides a benchmark dataset for this task. Most work uses a single model architecture
for both prediction and optimization. Our contribution is demonstrating that role
separation—matching architecture to task—outperforms monolithic deployment.

\textbf{Bayesian optimization for EDA.}
BO has been applied to VLSI placement~\cite{agnesina2020} and design-space
exploration~\cite{nardi2019}. We use qLogNEHVI~\cite{daulton2021,ament2024} for multi-objective
optimization with MC Dropout uncertainty, applied to a GAT surrogate trained on real
OpenROAD outputs. We evaluate RL performance using the Interquartile Mean (IQM)
metric~\cite{agarwal2021}, which is more robust than mean return for small seed counts.

\textbf{EDA reproducibility.}
OpenROAD~\cite{openroad2019} and ORFS provide open-source RTL-to-GDS flows.
We document that ORFS version changes produce $>$10\% PPA divergence on identical
inputs, including timing-sign reversals—a finding with implications for any ML-for-EDA
work using ORFS as ground truth. To our knowledge, no prior work has quantified
cross-version PPA sensitivity in open-source EDA flows.

%======================================================================
\section{Methodology}
\label{sec:method}
%======================================================================

\subsection{Pre-Registered Experimental Design}

Each phase had outcomes committed before execution:

\begin{itemize}
\item \textbf{Phase 1 (RL):} ``IQM $>$ 0 after 500K steps'' $\rightarrow$ FAIL
with monolithic GAT (IQM = $-$17.6). Root cause diagnosed via entropy ablation.
Threshold unchanged; architecture re-specified to GBDT $\rightarrow$ PASS
(IQM = +67.5)
\item \textbf{Phase 2a (GAT training):} ``LODO-CV MAE $<$ 0.2'' $\rightarrow$ PASS
(MAE = 0.153)
\item \textbf{Phase 2a (Criticality ablation):} ``$|\Delta\text{MAE}|>0.01$
$\rightarrow$ criticality helps; else redundant'' $\rightarrow$ NULL ($\Delta = -0.003$)
\item \textbf{Phase 2a (BO):} ``Pareto front grows by $\geq$1 candidate''
$\rightarrow$ PASS (+4 AES, +3 JPEG)
\item \textbf{Phase 2a (C+):} ``$\geq$4/7 Pareto-dominant $\rightarrow$
generalizes'' $\rightarrow$ PASS (7/7)
\end{itemize}

This structure prevents post-hoc metric selection and makes negative results
(criticality null, version sensitivity) first-class contributions.

\subsection{Surrogate Architecture and Role Separation}

The original framework deployed a GAT model for both RL reward evaluation and BO
surrogate prediction. We discovered this was architecturally mismatched:

\begin{itemize}
\item \textbf{GAT for RL:} The RL environment requires fast ($<$1ms), scalar reward
evaluation from tabular knob vectors. GAT requires graph construction per evaluation,
is slower, and its graph inputs were semantically mismatched to the RL action space
(single clock-target knob).

\item \textbf{GBDT for RL:} GBDT takes tabular features directly, evaluates in
$\sim$0.1ms, and achieves MAE = 0.153 on the same test set. Replacing GAT with
GBDT as the RL reward model produced IQM = +67.5 (vs.\ $-$17.6).

\item \textbf{GAT for BO:} GAT's graph-structured input captures netlist topology.
With MC Dropout (50 forward passes per candidate), it provides conservative
uncertainty estimates suitable for qLogNEHVI's exploration-favoring acquisition.
This is the task GAT was designed for.
\end{itemize}

\subsection{Data Pipeline}

\textbf{Training data.} 23 unique configurations across 4 ASAP7 RTL designs
(AES, Ibex, JPEG encoder, RISC-V 32I) at 5--7 clock targets per design,
synthesized through ORFS v3.0. Each run produces a placed-and-routed DEF file,
timing reports, and power/area metrics.

\textbf{Graph extraction.} DEF files are parsed into PyTorch Geometric graphs
(9K--62K nodes per design) with cell-type node features and net-connectivity edges.
100\% extraction success rate across all 23 configurations.

\textbf{Leave-One-Design-Out CV.} With only 4 design families, standard random
splits would leak design identity. LODO-CV holds out all configurations of one
design family, training on the other three.

\subsection{Sanity Checks as Methodology}

Every experimental phase began with a pre-flight check designed to catch systematic
errors before scaling. Four checks caught four distinct bugs:

\begin{enumerate}
\item \textbf{RL entropy ablation} $\rightarrow$ caught broken GAT predictor
(random-graph inputs producing noise-dominated rewards)
\item \textbf{Data provenance spot-check} $\rightarrow$ caught v1/v2 data confusion
(extraction initially ran on broken SDC-sweep data where all clock configs per family
produced byte-identical DEF files; MD5 check revealed only 4 unique designs instead of
23, making all downstream GAT predictions structurally constant per family)
\item \textbf{EDA calibration probe} $\rightarrow$ caught ORFS version mismatch
(v3.0 vs.\ 26Q1 producing $>$10\% PPA divergence)
\item \textbf{C+ implausible-value check} $\rightarrow$ caught SDC time-unit
mismatch (ASAP7 uses picoseconds; script specified nanoseconds)
\end{enumerate}

Each would have produced a plausible result if unchecked. The RL entropy ablation
would have concluded ``RL doesn't work for PPA optimization'' rather than ``the
predictor is broken.'' The v1/v2 data confusion would have trained a GAT on
graphs with zero within-family structural variation, producing a legitimate-looking
but meaningless criticality ablation. The SDC unit error would have reported
22$\times$ inflated power as a real EDA result. The ORFS version mismatch would
have produced verification verdicts that appeared interpretable but silently
compared surrogate predictions against a different EDA pipeline than training,
yielding pass/fail results that measured cross-version noise rather than
surrogate accuracy.

%======================================================================
\section{Experimental Results}
%======================================================================

\subsection{Finding 1: Role Separation (+67.5 IQM)}

The monolithic GAT-as-RL-surrogate produced IQM = $-$17.6 over 500K PPO steps
(5 seeds). Diagnosis revealed the GAT was receiving randomly-constructed graphs
with no semantic connection to the RL agent's clock-target actions.

Replacing the RL reward model with GBDT (keeping GAT for BO) produced
IQM = +67.5—an 85.1-point improvement. The RL agent learned a coherent
policy along the single dimension the GBDT surrogate could resolve (clock
target), confirming that role separation enabled training to proceed where
the monolithic architecture could not.

\subsection{Finding 2: GAT Surrogate Quality (MAE = 0.153)}

\begin{table}[t]
\centering
\caption{LODO-CV results on 23 real OpenROAD configurations}
\label{tab:lodo}
\begin{tabular}{lccc}
\toprule
\textbf{Hold-out} & \textbf{GAT+c(e)} & \textbf{GAT$-$c(e)} & \textbf{Size+clock LR} \\
\midrule
AES      & 0.175 & 0.118 & 0.322 \\
Ibex     & 0.100 & 0.108 & 0.071 \\
JPEG     & 0.271 & 0.272 & 0.679 \\
RISC-V   & 0.076 & 0.112 & 0.264 \\
\midrule
\textbf{Mean} & \textbf{0.156} & \textbf{0.153} & \textbf{0.376} \\
\bottomrule
\end{tabular}
\end{table}

Table~\ref{tab:lodo} shows LODO-CV results. The GAT without criticality edges
achieves mean MAE = 0.153, outperforming the size+clock linear regression baseline
(MAE = 0.376) by 2.5$\times$. This confirms the GAT learns genuine structural
features from netlist topology, not merely design size. The advantage is not
uniform—Ibex's small, regular structure is well-modeled by simple regression
(0.071 vs.\ 0.108), while JPEG's larger and more diverse netlist topology
benefits substantially from the GAT's structural representation
(0.272 vs.\ 0.679).

\subsection{Finding 3: Criticality Null with Mechanism}

Depth-based criticality edge weights produce $\Delta$MAE = $-$0.003 ($-$2.0\%)—
well within the pre-registered $\pm$0.01 indifference zone. Criticality edges
are redundant with information already encoded in graph topology.

\textbf{Mechanism.} Within each design family, the number of nodes in the
placed-and-routed netlist correlates perfectly with clock target
($\rho = -1.0$): tighter timing constraints trigger more buffer insertion by
the synthesis tool, increasing node count. The GAT already observes node count
through graph structure; adding explicit timing-criticality edges provides no
new information.

This is not a failure of the criticality concept—it is a finding about the
ASAP7/ORFS synthesis pipeline's buffer-insertion behavior at the single-knob
(clock target) optimization level. Multi-knob optimization (placement density,
aspect ratio, CTS parameters) may break this correlation, potentially
restoring criticality's value.

\subsection{Finding 4: Surrogate-Only Pareto Expansion}

\begin{table}[t]
\centering
\caption{Bayesian optimization results per design family (50 qLogNEHVI trials)}
\label{tab:bo}
\begin{tabular}{lccccc}
\toprule
\textbf{Family} & \textbf{Seed} & \textbf{Final} & \textbf{Growth} & \textbf{MC std} \\
\midrule
AES      & 5 & 9  & +4 & 0.080 \\
JPEG     & 6 & 9  & +3 & 0.089 \\
Ibex     & 6 & 6  & +0 & 0.094 \\
RISC-V   & 6 & 6  & +0 & 0.096 \\
\bottomrule
\end{tabular}
\end{table}

Table~\ref{tab:bo} summarizes BO outcomes. AES and JPEG show Pareto-front growth
(+4 and +3 new candidate configurations respectively), while Ibex and RISC-V
show no growth—the surrogate correctly identified that the existing manifests
already cover the achievable trade-off space for those designs.

\subsection{Finding 5: ORFS Version Sensitivity}

\begin{figure}[t]
\centering
\includegraphics[width=\columnwidth]{fig3_version.pdf}
\caption{PPA deltas between ORFS v3.0 and 26Q1 on identical RTL/SDC inputs.
The $\pm$5\% tolerance band is exceeded by all designs on at least one metric.
JPEG exhibits a WNS sign flip (timing met $\rightarrow$ timing violated).}
\label{fig:version}
\end{figure}

When we attempted EDA-in-the-loop verification, we discovered that ORFS v3.0
(used for training data) could not be built on current Ubuntu 22.04 due to
dependency rot. Using the available ORFS 26Q1, a calibration probe across 4
designs revealed erratic PPA divergence (Fig.~\ref{fig:version}).

The critical finding is directional inconsistency: power increases for AES
(+9.7\%) but decreases for JPEG ($-$6.5\%). WNS diverges by $-$97\% to
$-$193\% for three designs. Most concerning, JPEG's WNS flips sign entirely:
v3.0 reports timing met (+13.7ps slack); 26Q1 reports timing violated
($-$12.8ps).

This is not calibration error—it is two different synthesis/P\&R algorithms
making qualitatively different physical decisions on the same input. Any
ML-for-EDA work using ORFS as ground truth must control for version explicitly,
or risk training on one pipeline's decisions and evaluating against another's.

\subsection{Finding 6: Cross-Version Generalization}

\begin{figure}[t]
\centering
\includegraphics[width=\columnwidth]{fig2_tradeoff.pdf}
\caption{Trade-off direction for all 7 BO candidates relative to native-clock
baselines. All candidates reduce power at the cost of fmax, matching the
surrogate's predicted trade-off direction. AES candidates cluster at
$-$35\% to $-$45\% fmax; JPEG candidates spread across $-$27\% to $-$32\%.}
\label{fig:tradeoff}
\end{figure}

\begin{figure}[t]
\centering
\includegraphics[width=\columnwidth]{fig1_pareto.pdf}
\caption{Pareto fronts for AES and JPEG. BO candidates (stars) expand the
frontier toward lower power. CAND-6 strictly dominates CTRL-2 in the JPEG
family. Reference set: 26Q1 calibration baselines + C+ controls (2 points
per family).}
\label{fig:pareto}
\end{figure}

Despite Finding~5 establishing that v3.0 and 26Q1 produce systematically
different PPA values, we tested whether the surrogate's \textit{relative}
Pareto ordering survives cross-version transfer. Seven BO-proposed candidates
(4~AES, 3~JPEG) and 3 control configurations were synthesized through
26Q1.

\textbf{Trade-off direction is the primary result.} All 7 candidates
consistently expand the Pareto frontier toward lower power at the cost of
reduced fmax (Fig.~\ref{fig:tradeoff}), matching the surrogate's predicted
trade-off direction. This consistency across both design families and all 7
candidates is unlikely to occur by chance and suggests the surrogate learned
genuine design-space structure rather than version-specific EDA artifacts.

\textbf{Pareto status.} All 7 candidates are non-dominated against 26Q1
calibration baselines and validation controls (2 reference points per design
family, Fig.~\ref{fig:pareto}). The small reference set limits the strength
of this claim. However, one candidate (CAND-6) strictly dominates control
CTRL-2 in the JPEG family—better power (0.0683W vs.\ 0.0701W), smaller area
(6386 vs.\ 6389 $\mu$m$^2$), and higher fmax (1262.75 vs.\ 1217.56 MHz).
This demonstrates that BO is capable of surfacing genuinely Pareto-improving
configurations, not merely non-dominated ones in a sparse space. (The area
difference of 3~$\mu$m$^2$ is small in isolation, but dominance holds even
if area is treated as tied: the power and fmax improvements each exceed
the respective metric's resolution.)

All 7 candidates and 3 controls met timing (WNS = 0.00ps), confirming that
the BO pipeline produces physically realizable designs.

%======================================================================
\section{Compute Cost}
%======================================================================

\begin{figure}[t]
\centering
\includegraphics[width=0.85\columnwidth]{fig4_cost.pdf}
\caption{Per-phase compute cost breakdown. Total: \$8.35 on AWS (spot +
on-demand c6i.4xlarge instances). The diagnostic phase that discovered ORFS
version sensitivity cost \$0.65—the cheapest phase, yet arguably the most
impactful finding.}
\label{fig:cost}
\end{figure}

The full experimental arc cost \$8.35 (Fig.~\ref{fig:cost}): \$0.85 for
bootstrap and qLogNEHVI migration (t3.medium spot, $\sim$2h), \$1.20 for
RL training with 500K PPO steps across 5 seeds (t3.medium spot, $\sim$3h),
\$2.95 for real-data graph extraction, GAT retraining, and BO
(c6i.xlarge spot instances across multiple sessions, $\sim$6h total),
\$0.65 for the calibration probe that discovered version sensitivity
(c6i.4xlarge on-demand, $\sim$1h), and \$2.70 for the C+ generalization
probe including failed first batch and re-runs (c6i.4xlarge on-demand,
$\sim$4h).

All computation ran on AWS c6i.4xlarge instances (16 vCPU, 32 GB RAM,
\$0.68/hr on-demand). No GPUs were used. EDA run times ranged from 8--9
minutes (AES) to 48--55 minutes (JPEG) per configuration on 26Q1.

We report cost explicitly because most ML-for-EDA papers do not, making
it difficult for other researchers to assess whether results are achievable
at academic budgets. Our total spend of \$8.35 for the complete experimental
arc—training, optimization, ablation, and EDA verification—is comparable to
a few GPU-hours on cloud instances, demonstrating that meaningful ML-for-EDA
experimentation is feasible without dedicated hardware.

%======================================================================
\section{Discussion}
%======================================================================

\subsection{The Methodology Is the Contribution}

The six findings span positive results (role separation, Pareto expansion,
cross-version generalization), a mechanistically-explained null (criticality
redundancy), and a reproducibility hazard (ORFS version sensitivity). No
single finding would carry a paper on its own. Together, they demonstrate
what a rigorously instrumented ML-for-EDA pipeline looks like in practice.

The pre-registration discipline is load-bearing, not decorative. Without it:
\begin{itemize}
\item The criticality null would likely have been omitted (unpublishable
negative result)
\item The version sensitivity finding would have been an undocumented
debugging session rather than a contribution
\item The C+ probe would have been framed as ``verification'' rather than
honestly labeled as a generalization experiment
\end{itemize}

\subsection{Limitations}

\textbf{Scale.} 4 designs, 23 configurations, single-knob (clock target)
optimization. Multi-knob, multi-design-family experiments are needed to test
whether the criticality null persists and whether cross-version generalization
holds at scale.

\textbf{Reference set size.} The C+ Pareto analysis uses 2 reference points
per design family. While CAND-6's strict domination is unambiguous, the
remaining 6 candidates' non-dominated status would be more convincing against
a larger reference set.

\textbf{No true verification.} The v3.0 build is unrecoverable on current
systems, preventing direct comparison between surrogate predictions and the
EDA pipeline they were trained on. The C+ result is a cross-version
generalization claim, not a verification claim.

\textbf{Surrogate limitations.} MC Dropout overestimates uncertainty (U/S
ratio 6--9$\times$ per objective), causing conservative BO exploration. This
is the favorable failure mode: over-estimated uncertainty drives conservative
acquisition, making the Pareto-expansion claims in Finding~4 more robust than
they would be under a well-calibrated or under-confident surrogate. Better
calibrated uncertainty (e.g., deep ensembles) could improve Pareto
expansion further.

\subsection{Future Directions}

Three extensions follow directly from the current limitations.
First, the criticality null (Finding~3) is specific to single-knob
optimization where buffer insertion creates a perfect correlation between
clock target and node count; multi-knob optimization (placement density,
aspect ratio, CTS parameters) may break this correlation and restore
criticality's predictive value.
Second, the ORFS version sensitivity finding motivates a community resource:
a versioned-ORFS reference benchmark with frozen PPA baselines, enabling
cross-study comparison without version-confounded results.
Third, the C+ generalization result—7/7 surrogate-predicted Pareto candidates
confirmed on a different ORFS version—motivates systematic testing across
additional ORFS releases and PDKs beyond ASAP7 to establish how far
surrogate transferability extends.

%======================================================================
\section{Conclusion}
%======================================================================

We presented an ML-driven PPA optimization framework distinguished by its
methodological instrumentation: pre-registered outcomes, systematic sanity
checks, and honest reporting of negative results. The framework produced six
findings across the experimental arc, including a surrogate role-separation
insight, a mechanistically-explained criticality null, and evidence that
Bayesian optimization's Pareto-ordering generalizes across ORFS versions
despite $>$10\% absolute PPA divergence.

The ORFS version sensitivity finding (Section~IV-E) has implications beyond
this project: any ML-for-EDA pipeline using ORFS as ground truth must lock
and report the exact ORFS version, as version changes can flip timing
verdicts on identical inputs.

The total compute cost of \$8.35 demonstrates that meaningful ML-for-EDA
experimentation—including real OpenROAD synthesis, EDA verification, and
multiple ablation controls—is feasible for individual researchers on
commodity cloud hardware without GPUs or institutional infrastructure.

\subsection*{Reproducibility}

As an independent research project, the complete experimental record is the
primary deliverable. All code, data manifests, trained models, ORFS run
reports, and pre-registered plan documents are available at
\url{https://github.com/oguzhan-canada/instrumented-ml-ppa}.
The experimental trail (PROGRESS.md) documents every decision, bug fix,
and sanity-check result in chronological order.
Artifacts are organized by phase: Phase~0 bootstrap scripts and qLogNEHVI
migration; Phase~1 RL training checkpoints, entropy ablation logs, and IQM
curves; Phase~2a GAT weights, LODO-CV splits, BO trial histories, and
criticality ablation results; and C+ probe EDA verification outputs
including full ORFS timing/power/area reports for all 10 candidate and
control configurations.
Pre-registered experimental plans with success/failure thresholds are
committed before each phase's execution, enabling independent audit of
the pre-registration claims made in Section~\ref{sec:method}.

\section*{Acknowledgments}

This work was conducted as independent research by the sole author,
without institutional or external funding. All compute was self-funded
on AWS commodity instances. The author thanks the OpenROAD, ASAP7,
BoTorch, and Stable-Baselines3 open-source communities, whose tools
made this project feasible at the reported budget.

%======================================================================
% References
%======================================================================
\bibliographystyle{IEEEtran}
% Placeholder references — replace with .bib file
\begin{thebibliography}{13}

\bibitem{kahng2021}
A.~B.~Kahng, ``Machine learning applications in physical design,''
\textit{Proc. Int. Symp. Physical Design (ISPD)}, pp.~68--75, 2021.

\bibitem{openroad2019}
T.~Ajayi \textit{et al.}, ``INVITED: toward an open-source digital flow:
first learnings from the OpenROAD project,''
\textit{Proc. DAC}, pp.~1--4, 2019.

\bibitem{cheng2021}
W.-H.~Cheng \textit{et al.}, ``On joint learning for solving placement and
routing in chip design,'' \textit{Proc. NeurIPS}, 2021.

\bibitem{huang2021}
G.~Huang, J.~Hu, Y.~He, J.~Liu, M.~Ma, Z.~Shen, J.~Wu, Y.~Xu,
H.~Zhang, K.~Zhong, X.~Ning, Y.~Dong, Y.~Wang, H.~Yang, and B.~Yu,
``Machine learning for electronic design automation: a survey,''
\textit{ACM Trans. Des. Autom. Electron. Syst.},
vol.~26, no.~5, art.~40, pp.~1--46, 2021.

\bibitem{wu2022}
Y.-C.~Wu \textit{et al.}, ``High-definition routing congestion prediction
for large-scale FPGAs,'' \textit{Proc. ASPDAC}, 2022.

\bibitem{ustun2020}
E.~Ustun, C.~Deng, D.~Pal, Z.~Li, and Z.~Zhang,
``Accurate operation delay prediction for FPGA HLS using graph neural
networks,'' \textit{Proc. ICCAD}, pp.~1--9, 2020.

\bibitem{chhabria2022}
V.~A.~Chhabria \textit{et al.}, ``Towards machine learning for placement
and routing in chip design: a methodological overview,''
\textit{arXiv:2202.13564}, 2022.

\bibitem{chowdhury2021}
A.~B.~Chowdhury \textit{et al.}, ``OpenABC-D: a large-scale dataset for
machine learning guided integrated circuit synthesis,''
\textit{Proc. DATE}, pp.~1696--1701, 2021.

\bibitem{agnesina2020}
A.~Agnesina, P.~Rajvanshi, T.~Yang, G.~Pradipta, and S.~K.~Lim,
``VLSI placement parameter optimization using deep reinforcement learning,''
\textit{Proc. ICCAD}, pp.~1--9, 2020.

\bibitem{nardi2019}
L.~Nardi, D.~Koeplinger, and K.~Olukotun,
``Practical design space exploration,''
\textit{Proc. MASCOTS}, pp.~347--358, 2019.

\bibitem{daulton2021}
S.~Daulton, M.~Balandat, and E.~Bakshy,
``Parallel Bayesian optimization of multiple noisy objectives with expected
hypervolume improvement,''
\textit{Proc. NeurIPS}, vol.~34, 2021.

\bibitem{ament2024}
S.~Ament, S.~Daulton, D.~Eriksson, M.~Balandat, and E.~Bakshy,
``Unexpected improvements to expected improvement for Bayesian optimization,''
\textit{Proc. NeurIPS}, vol.~36, 2023.

\bibitem{agarwal2021}
R.~Agarwal, M.~Schwarzer, P.~S.~Castro, A.~Courville, and M.~G.~Bellemare,
``Deep reinforcement learning at the edge of the statistical precipice,''
\textit{Proc. NeurIPS}, 2021.

\end{thebibliography}

\balance
\end{document}
